\begin{lemma}\label{lem:decic77-geometry}
Put $X=X_1(77)$ and
$Y=X/\langle\langle36\rangle\rangle$.
Then $Y$ has good reduction at $2$, genus $37$, and
$T_2J_Y(\mathbb Q)$ is finite.
The cover $Y\to X_0(77)$ is etale cyclic of degree $6$.
Its deck group is generated by $\delta=\langle2\rangle$, and
$\langle3\rangle=\delta^{-1}$.
\end{lemma}
\begin{proof}
Let $\pi:X\to Y$ be the quotient by
$\langle\langle36\rangle\rangle$, of order $5$.
This group acts freely, including on the special fibre at $2$.
Thus $Y$ has good reduction at $2$, and
$g(Y)=1+(181-1)/5=37$.
The complete degree-independent winding computation accounts
for all $181$ source dimensions.
Every nontrivial-character component is full winding.
The trivial-character component has dimension $7$ and winding
dimension $6$, and $T_2$ is zero on its one-dimensional quotient.
Kato's theorem, in the winding formulation of
\cite[Remarks~3.2 and 3.3]{DEHMZ}, therefore gives that $T_2J_X(\mathbb Q)$ is finite.
Pullback and norm show that $T_2J_Y(\mathbb Q)$ is finite as well.

We retain the exact modular identification of $Y_{\mathbb F_2}$.
The base $X_0(77)$ is the genus-$7$ normalization obtained from
$X_0(11)$ by adjoining $v$ with
\[
 v^8+jv^7+v^4+1=0.
\]
Its Picard group is $C_{30}\times C_{60}\times\mathbb Z$.
The degree-$6$ cyclic cover has cusp Frobenius elements
$0,2,3,5$ at widths $77,11,7,1$.
The cusp conditions leave six characters to $C_6$; splitting
at two actual degree-$10$ source images singles out the modular
character.
Its cubic subcover is
$X_1(7)\times_{X(1)}X_0(11)$, with Tate parameter $t$ and
\[
 b=t^3-t^2,\quad c=t^2-t,\quad
 v=\frac{t^2+t}{t^3+t+1}.
\]
The fixed Artin--Schreier component is checked on generators of
the base Picard group modulo $2$.
The composite has genus $37$, a rational cusp and $j$-degree
$576$. Its two physical deck transformations are rebuilt and
verified exactly.


Since $2^6=36^2$ and no smaller positive power of $2$ belongs
to $\{\pm36^i:0\leq i<5\}$, $\delta$ has order $6$.
Also $36^3=71\equiv-6\pmod{77}$, giving
$\langle3\rangle\delta=1$.
In the explicit cubic and Artin--Schreier model, with coordinates
$(t,s,z)$, its physical action is
\[
 \delta(t,s,z)=\left(\frac{t+1}{t},s,z+1\right).
\]
The Tate normal form of $2P_7$ verifies the first coordinate;
the quadratic character has value $-1$ at $2$.
Both component actions and the order-$6$ identity are checked
exactly on the defining function fields.
\end{proof}
